\section{Discussion}
\label{sec:extensions}

\paragraph{Effect-ordering (mixed-integer).}
When the restriction set is the effect-ordering polyhedron of Restriction~1 of \citet{boyarsky2026}, the program \eqref{eq:surrogate} becomes a mixed-integer
LP. The locally-constant integer-solution argument of \S1(ii) in \citet{boyarsky2026} extends to this case. On a neighborhood of $\gamma_\mathcal{S}$, the
optimal sign pattern is constant, so the MILP reduces to the LP
\eqref{eq:surrogate} after fixing integer variables. We then perform selection
by \eqref{eq:surrogate} at the fixed pattern. Sign flips across candidate sites
can be detected by re-solving the MILP at the top two or three candidates.

\paragraph{Cost-weighted selection.}
Adding a cost $c_k$ to the objective \eqref{eq:minimax} is straightforward
because the cost term is separable in $k$. The minimax-optimal candidate under
Lagrangian \eqref{eq:lagrangian} is obtained by solving \eqref{eq:surrogate}
with an additional additive penalty. In our setting a cost based on ethical
concerns might penalize candidates that are demographically similar to existing
sites, because such sites do little to broaden the external-validity claim
despite being cheap to reach.

\paragraph{Multiple candidates at once.}
The design problem \eqref{eq:minimax} can be posed over a subset
$\mathcal{J} \subseteq \mathcal{K}$ of candidates to add simultaneously; in this
case the width surrogate \eqref{eq:surrogate} becomes a set-function in
$\mathcal{J}$ that is \emph{submodular} in the rank-monotone limit. That is,
adding the first hull-expander contributes more than adding the second. A greedy
algorithm attains a $(1 - 1/e)$-approximation under mild regularity, which
matches the classical D-optimality result of \citet{nemhauser1978analysis}.

\paragraph{Connection to efficient influence functions.}
The minimum-norm multiplier $\lambda^\star$ that appears in \eqref{eq:envelope}
is the same object as the Riesz representer in the efficient influence function of \citet[Eq.~2.9]{boyarsky2026}. The minimax design criterion therefore gets
its inference validity from their debiased cross-fitting. Confidence intervals
for the post-addition width are a standard delta-method computation.

\paragraph{Open questions.}
Three directions are natural follow-ups. First, the adversarial $b$-prior of
Section~\ref{sec:unknown-b}(3) admits a saddle-point representation but not a
closed-form solution. Existing interior-point solvers handle small instances,
but large candidate pools would benefit from a custom Frank--Wolfe scheme.
Second, in the online-experimentation setting the candidate pool is revealed
sequentially, and the selection must be made without access to the full set.
Adapting \eqref{eq:surrogate} to this regret-minimization formulation is open.
Third, the assumption that the within-site $X$-distribution at candidate sites
is known (or has an informative prior) is strong. A distributionally-robust
variant that treats $P_X^{(k)}$ as uncertain over a Wasserstein ball would keep
the theoretical guarantees of the \citet{jeong2020assessing} worst-case program,
at the cost of a more conservative selection.
